\documentclass[11pt]{article}
\usepackage{amsmath,amssymb,amsthm}
\usepackage[margin=1in]{geometry}

\title{Disproof of a Claim About $\mathbb{E}[\hat N]$}
\date{}

\newtheorem{proposition}{Proposition}
\newtheorem{remark}{Remark}

\begin{document}
\maketitle

\paragraph{Context (Sheet 02, Problem 2).}
In the capture--recapture simulation used in \texttt{sheet-02/code.ipynb}, the estimator of the population size is
\[
\hat N \;=\; n + \hat f_0,
\]
where $f_k$ denotes the number of individuals caught exactly $k$ times, $n=\sum_{k\ge 1} f_k$ is the number of distinct observed individuals, and
\[
\hat f_0 \;=\;
\begin{cases}
\dfrac{f_1(f_1-1)}{2(f_2+1)}, & f_2=0,\\[6pt]
\dfrac{f_1^2}{2f_2}, & f_2>0.
\end{cases}
\]

\paragraph{Claim to assess.}
The statement in question reads (rewritten with standard expectation notation):
\begin{quote}
``$\mathbb{E}[\hat N] > N$, or if $\mathbb{E}[\hat N]=N$, then $\mathbb{E}[\hat N-N] > 0$ for $n<\infty$.'' 
\end{quote}

\begin{proposition}[The second clause is impossible]
For any integrable estimator $\hat N$ of a fixed (non-random) parameter $N$,
\[
\mathbb{E}[\hat N-N] \;=\; \mathbb{E}[\hat N] - N.
\]
In particular, if $\mathbb{E}[\hat N]=N$, then $\mathbb{E}[\hat N-N]=0$.
\end{proposition}

\begin{proof}
Linearity of expectation yields
\[
\mathbb{E}[\hat N-N]=\mathbb{E}[\hat N]-\mathbb{E}[N]=\mathbb{E}[\hat N]-N,
\]
because $N$ is a constant. The special case follows immediately.
\end{proof}

\begin{remark}
So the implication ``if $\mathbb{E}[\hat N]=N$ then $\mathbb{E}[\hat N-N]>0$'' is false for all finite or infinite sample sizes; it contradicts basic properties of expectation.
\end{remark}

\begin{proposition}[The strict inequality $\mathbb{E}(\hat N)>N$ is not generally true]
There is no general principle forcing $\mathbb{E}[\hat N]>N$ for estimators arising in capture models at finite sample size.
\end{proposition}

\begin{proof}
Consider the natural estimator $\tilde N := n$ (the number of distinct observed individuals), which is a standard baseline in the same capture setup.

Assume a simple model consistent with the simulation spirit: there are $N$ individuals, and across $T$ capture occasions each individual is caught independently with per-occasion probability $p\in(0,1)$. Then an individual is observed at least once with probability
\[
q \;=\; 1-(1-p)^T \in (0,1)\qquad (T<\infty),
\]
and therefore
\[
n \sim \mathrm{Binomial}(N,q)
\quad\Rightarrow\quad
\mathbb{E}[\tilde N]=\mathbb{E}[n]=Nq < N.
\]
Hence even in a basic finite-sampling capture model, one can have $\mathbb{E}[\hat N]<N$ for a reasonable estimator. This disproves any unconditional claim that ``the expected value must exceed $N$.'' 
\end{proof}

\paragraph{Conclusion.}
The statement is \emph{false as written}. The second clause is mathematically impossible, and the first clause is not a general truth (even within simple capture models for finite sampling).

\paragraph{How to proceed if you know the catch probability distribution (e.g.\ $p_i\sim U(0,a)$).}
If your goal is instead to study a \emph{one-sided probability} like
\[
\Pr(\hat N > N) > \Pr(\hat N \le N)
\qquad\Longleftrightarrow\qquad
\Pr(\hat N > N)>\tfrac12,
\]
then you need distributional information about $\hat N$. In the capture setup of the notebook, a common route is:

\subparagraph{Step 1: Specify the per-individual catch-count model.}
Fix the number of capture occasions $T<\infty$. For individual $i\in\{1,\dots,N\}$ let
\[
K_i \mid p_i \sim \mathrm{Binomial}(T,p_i),
\]
where $p_i\in(0,1)$ is that individual's per-occasion catch probability. Assume $(p_i)_{i=1}^N$ are i.i.d.\ from some distribution $G$ (e.g.\ $U(0,a)$).

\subparagraph{Step 2: Compute the mixed probabilities $q_k=\Pr(K_i=k)$.}
Unconditionally, the $K_i$ are i.i.d.\ with the \emph{mixture} pmf
\[
q_k \;:=\; \Pr(K_i=k)
\;=\;
\binom{T}{k}\,\mathbb{E}_{p\sim G}\!\left[p^k(1-p)^{T-k}\right],
\qquad k=0,1,\dots,T.
\]
For $p\sim U(0,a)$ one gets an explicit integral:
\[
q_k
=
\binom{T}{k}\frac{1}{a}\int_0^a p^k(1-p)^{T-k}\,\mathrm dp
=
\binom{T}{k}\frac{B_a(k+1,T-k+1)}{a},
\]
where $B_a(\alpha,\beta)=\int_0^a t^{\alpha-1}(1-t)^{\beta-1}\,\mathrm dt$ is the (lower) incomplete beta function.

\subparagraph{Step 3: Use the multinomial law for the frequency counts.}
Define $f_k:=\sum_{i=1}^N \mathbf{1}\{K_i=k\}$. Since the $K_i$ are i.i.d.\ with category probabilities $(q_0,\dots,q_T)$,
\[
(f_0,f_1,\dots,f_T)\sim \mathrm{Multinomial}\!\left(N;\,q_0,q_1,\dots,q_T\right).
\]
In particular,
\[
\mathbb{E}[f_k]=Nq_k,\quad
\mathrm{Var}(f_k)=Nq_k(1-q_k),\quad
\mathrm{Cov}(f_k,f_\ell)=-Nq_kq_\ell\ (k\neq \ell).
\]
Also $n=\sum_{k\ge 1} f_k = N-f_0$, hence $\mathbb{E}[n]=N(1-q_0)$ and $\mathrm{Var}(n)=Nq_0(1-q_0)$.

\subparagraph{Step 4: Approximate the distribution of $\hat N$ (delta method).}
Ignoring the $f_2=0$ fallback for the moment (see Step 5), the estimator used in the notebook is
\[
\hat N \;=\; g(n,f_1,f_2)
\;:=\;
n+\frac{f_1^2}{2f_2}\qquad(f_2>0).
\]
For large $N$, the multinomial central limit theorem implies
\[
(n,f_1,f_2)\approx \mathcal N(m,\Sigma),
\quad
m=\big(N(1-q_0),\,Nq_1,\,Nq_2\big),
\]
with $\Sigma$ determined by the multinomial covariances. The delta method then gives an approximate normal law
\[
\hat N \approx \mathcal N\!\big(\mu_{\hat N},\,\sigma_{\hat N}^2\big),
\qquad
\mu_{\hat N}:=g(m)=N\Big(1-q_0+\frac{q_1^2}{2q_2}\Big),
\]
and
\[
\sigma_{\hat N}^2 \approx \nabla g(m)^\top \Sigma\,\nabla g(m),
\quad
\nabla g(m)=\left(1,\ \frac{q_1}{q_2},\ -\frac{q_1^2}{2q_2^2}\right).
\]
Using $\mathrm{Cov}(n,f_1)=Nq_0q_1$, $\mathrm{Cov}(n,f_2)=Nq_0q_2$, $\mathrm{Cov}(f_1,f_2)=-Nq_1q_2$, you can write $\sigma_{\hat N}^2$ explicitly as
\[
\frac{\sigma_{\hat N}^2}{N}\approx
q_0(1-q_0)
\;+\;\Big(\frac{q_1}{q_2}\Big)^2 q_1(1-q_1)
\;+\;\Big(\frac{q_1^2}{2q_2^2}\Big)^2 q_2(1-q_2)
\;+\;2\Big(\frac{q_1}{q_2}\Big)q_0q_1
\;-\;2\Big(\frac{q_1^2}{2q_2^2}\Big)q_0q_2
\;-\;2\Big(\frac{q_1}{q_2}\Big)\Big(\frac{q_1^2}{2q_2^2}\Big)q_1q_2.
\]

\subparagraph{Step 5: Compute $\Pr(\hat N>N)$ and check if it exceeds $1/2$.}
Under the normal approximation,
\[
\Pr(\hat N > N)\approx 1-\Phi\!\left(\frac{N-\mu_{\hat N}}{\sigma_{\hat N}}\right).
\]
Therefore \(\Pr(\hat N>N)>\tfrac12\) holds approximately whenever \(\mu_{\hat N}>N\), i.e.
\[
1-q_0+\frac{q_1^2}{2q_2} > 1
\quad\Longleftrightarrow\quad
\frac{q_1^2}{2q_2} > q_0,
\]
and the inequality becomes more decisive the larger \((\mu_{\hat N}-N)/\sigma_{\hat N}\) is.

\subparagraph{Step 6: Account for the $f_2=0$ fallback (if needed).}
In the multinomial model, $f_2\sim \mathrm{Binomial}(N,q_2)$ marginally, so
\[
\Pr(f_2=0)=(1-q_2)^N.
\]
If $Nq_2$ is moderately large, this probability is negligible and the $f_2>0$ approximation dominates. Otherwise, you must include the fallback definition
\(\hat f_0=\frac{f_1(f_1-1)}{2(f_2+1)}\) at $f_2=0$ and compute (analytically or by Monte Carlo) the mixture probability
\[
\Pr(\hat N>N)=\Pr(\hat N>N\mid f_2=0)\Pr(f_2=0)+\Pr(\hat N>N\mid f_2>0)\Pr(f_2>0).
\]

\end{document}

